GPU 版の OUxSBLI\_mini は，CUDA Fortran を用いて GPU 上で計算を実行する．本セクションでは GPU 版独有の実装について解説する．

\subsection{GPU メモリ確保と初期化}

GPU では明示的にメモリを確保する必要がある：

\begin{lstlisting}[language=Fortran]
real(8), allocatable, device :: Q(:,:,:,:)
allocate(Q(5,nx,ny,nz))
call check_gpu_error(cudaMemcpy(...))
\end{lstlisting}

CPU から GPU への初期データ転送も明示的に行う．

\subsection{GPU カーネルの構造}

CUDA Fortran のカーネルは \texttt{attributes(global)} で修飾される：

\begin{lstlisting}[language=Fortran]
attributes(global) subroutine calc_keep_x6(nx, ny, nz, Q, T, E)
  integer, intent(in), value :: nx, ny, nz
  real(8), intent(in), device :: Q(5,nx,ny,nz), T(nx,ny,nz)
  real(8), intent(out), device :: E(5,nx-1,ny-2,nz-2)

  integer :: i, j, k
  i = (blockIdx%x - 1) * blockDim%x + threadIdx%x
  j = (blockIdx%y - 1) * blockDim%y + threadIdx%y
  k = (blockIdx%z - 1) * blockDim%z + threadIdx%z

  if (i < nx .and. j < ny .and. k < nz) then
    ! thread ごとの計算
  endif
end subroutine
\end{lstlisting}

\subsection{Shared Memory の活用}

Shared memory は L1 キャッシュのように機能し，高速アクセスが可能である．OUxSBLI では KEEP flux の計算で shared memory を多用している：

\begin{lstlisting}[language=Fortran]
attributes(global) subroutine calc_keep_x6(...)
  real(8), dimension(:), shared :: rho, u, v, w, p, tmp
  integer :: ii, sx, sy, sz

  ! Global memory から shared memory に読み込み
  do ii = threadIdx%x - 2, threadsE%x + 3, blockDim%x
    ! thread が協調してデータを読み込み
    rho(ii) = Q(1,i+ii-threadIdx%x, j, k)
  enddo

  call syncthreads()  ! 全 thread の同期を待つ

  ! Shared memory からのアクセスは高速
  E(1,i,j,k) = KEEP6_x(rho, u, v, w, p, tmp)
end subroutine
\end{lstlisting}

Shared memory の利点：

\begin{enumerate}
  \item \textbf{低遅延}: 主メモリアクセスより $\sim 10$ 倍高速
  \item \textbf{同期機構}: 複数スレッド間のデータ交換が容易
  \item \textbf{メモリバンド幅}: L1 キャッシュと同等の高いバンド幅
\end{enumerate}

制限事項：

\begin{enumerate}
  \item \textbf{容量制限}: GPU 当たり数十 KB 程度（GPU 世代に依存）
  \item \textbf{スレッドブロック内のみ}: 異なるブロック間での共有不可
  \item \textbf{明示的な同期が必要}: syncthreads() での全スレッド待機
\end{enumerate}

\subsection{MPI による計算と出力の分離}

OUxSBLI\_mini では計算と出力を別々の MPI プロセスで実行する：

\begin{lstlisting}[language=Fortran]
if (mod(myrank, 2) == 0) then
  ! ランク 0, 2, 4, ... で計算
  call RungeKutta(myrank, mygpu, nx, ny, nz, ...)
else
  ! ランク 1, 3, 5, ... で出力
  call send_recv_for_print_odd(myrank, nranks, ...)
endif
\end{lstlisting}

利点：

\begin{enumerate}
  \item \textbf{計算と出力を並列化}: I/O 時間が計算時間に隠蔽される
  \item \textbf{スケーラビリティ}: 複数 GPU で効率的に計算
\end{enumerate}

通信オーバーヘッド：

\begin{enumerate}
  \item MPI\_ISEND/MPI\_IRECV で非同期通信
  \item 計算の合間に出力リクエストを処理
\end{enumerate}

\subsection{GPU カーネル呼び出し}

GPU 上のカーネル実行は以下の構文で呼び出される：

\begin{lstlisting}[language=Fortran]
call calc_keep_x<<<blocksE, threadsE, 1>>>(id_accuracy, nx, ny, nz, Q, T, E, sigma)
\end{lstlisting}

ここで：

\begin{itemize}
  \item \texttt{blocksE}: ブロック数（3D）
  \item \texttt{threadsE}: スレッド数（3D，通常は 32×4×1 など）
  \item \texttt{1}: Shared memory 使用量（バイト）
\end{itemize}

GPU では自動的に：

\begin{enumerate}
  \item スレッド数 $\times$ ブロック数個の並列スレッドを起動
  \item 各スレッド内の \texttt{threadIdx}, \texttt{blockIdx} を自動設定
  \item 全スレッド完了まで待機
\end{enumerate}

\subsection{GPU での精度特性}

GPU 版では 64 ビット（倍精度）浮動小数点演算を使用． CPU 版と同じ精度で実行される．

\end{document}
